\begin{table}[htbp]\centering
\begin{threeparttable}
\caption{The ITS rule and the species complexes moved to the same percentile.}
\label{tab:percentile}
\footnotesize\setlength{\tabcolsep}{4pt}
\begin{tabular}{lrrrrrr}
\toprule
Percentile & ITS limit & Complexes & Known demoted & False known, & False known, & Outside\\
 & (median) & (species) & (of 755) & species holdout (of 803) & V18 removed (of 673) & a complex\\
\midrule
95th & 4.4\% & 3 (6) & 32 (4.2\%) & 9 (1.1\%) & 37 (5.5\%) & 5\\
97.5th & 5.0\% & 3 (7) & 22 (2.9\%) & 13 (1.6\%) & 53 (7.9\%) & 18\\
99th & 6.0\% & 4 (9) & 6 (0.8\%) & 18 (2.2\%) & 53 (7.9\%) & 0\\
\bottomrule
\end{tabular}
\begin{tablenotes}\footnotesize
\item The ITS rule is re-applied to the stored decision-tree rows at each percentile, with the limit fitted
on each fold's training cultures, and the complexes are those of the distance rule at the same percentile
(\texttt{percentile\_sensitivity.py}). \emph{Known demoted}: known-species copies the rule demotes in the
culture holdout. \emph{V18 removed}: first-outcome M4 folds. The 99th-percentile row reproduces the reported
results exactly. Every error outside a complex is \emph{R.\ irregularis} called \emph{{[}R.{]}\ vesiculiferum}.
\end{tablenotes}
\end{threeparttable}
\end{table}
